% Volume 11: Design
% Opener: tagline + scope. Sits immediately after the \part{} declaration.

\vspace*{1.5cm}
\begin{center}
\itshape\large
Breakage becomes discipline
\end{center}
\vspace{0.5cm}

\noindent
Volumes I through X gave the reader the tools (Quantity, Form, Force), the materials (Matter), the substrate of action (Energy, Life), the languages of structure (Information, Machines, Systems), and the literacy of failure (Failure). This volume is the practice that integrates them: design. Design is what engineering is when seen from the engineer's chair. The reader who has worked through the previous ten volumes has the components from which design proceeds; this volume puts the components to work.

The volume's central claim is that design is the synthesis of constraint, want, knowledge, and judgment under risk. The four elements are distinct: constraint is what physics and budget and regulation will permit; want is what the user, the client, the community, and the future need; knowledge is what the engineer brings (from training, from experience, from the literature, from collaborators); judgment is what the engineer commits when the previous three under-determine the answer. Every chapter of this volume develops one or more facets of the synthesis.

The arc of the chapters: Chapter 1 defines design as a discipline. Chapter 2 covers requirements and specification, the discipline by which want becomes engineering input. Chapter 3 covers user research and human-centred design, the discipline by which the want is verified rather than assumed. Chapter 4 covers conceptual design and sketching, the discipline of generating and selecting candidate solutions. Chapter 5 covers detailed design and analysis, the discipline of turning a concept into something that can be built. Chapter 6 covers simulation and digital twins, the discipline of testing the design before it exists physically. Chapter 7 covers prototyping and iteration, the discipline of learning from things-that-do-not-yet-work. Chapter 8 covers testing and validation, the discipline of evidence that the design meets its requirements. Chapter 9 covers design reviews, certification, and sign-off, the institutional discipline of accountability. Chapter 10 covers manufacturability, maintainability, sustainability, and accessibility, the four "ilities" that determine whether the design serves its full lifecycle. Chapter 11 covers ethics under risk and uncertainty, treated here as a constitutive feature of design rather than a postscript. Chapter 12 is the integrated design studio: extended case studies done at full depth. Chapter 13 is the capstone, where the reader takes the entire book and applies it to one problem of meaningful scale.

The reader who finishes this volume should be able to take an engineering problem from need to deliverable design pack, with the requirements documented, the concepts compared, the detailed design analysed, the simulations conducted, the prototypes tested, the validation evidence assembled, the reviews satisfied, the manufacturability assessed, the lifecycle considered, and the ethics articulated. The deliverable is not academic exercise but professional output: a design pack a competent design panel would accept.

The volume's archetypes are predominantly ethics, optimisation, and interface; with substantial use of failure (in the design-against-failure sense developed in Volume X) and balance (in the tradeoff-discipline sense). The cases draw from across the engineering disciplines, with the discipline of design itself as the integrating thread.

\vspace{1cm}
